\section{Objectifs}
\begin{itemize}
\item Modéliser une croissance linéaire.
\item Utiliser une suite arithmétique pour une évolution discrète.
\item Utiliser une fonction affine pour une évolution continue.
\item Résoudre un problème de seuil.
\end{itemize}

\section{1. Croissance linéaire}
Une grandeur évolue linéairement lorsque son \textbf{accroissement} est constant.

Si un budget augmente de 120 € chaque mois et si $u(n)$ désigne sa valeur après $n$ mois :
\[
u(n+1)=u(n)+120.
\]

\section{2. Suites arithmétiques}
Une suite arithmétique de raison $r$ vérifie
\[
u(n+1)=u(n)+r.
\]

Si le terme initial est $u(0)$ :
\[
u(n)=u(0)+nr.
\]

\begin{itemize}
\item si $r>0$, la suite est croissante ;
\item si $r<0$, elle est décroissante ;
\item si $r=0$, elle est constante.
\end{itemize}

\section{3. Fonctions affines}
Une fonction affine s'écrit
\[
f(x)=ax+b.
\]

Le coefficient $a$ représente un taux d'accroissement constant : lorsque $x$ augmente d'une unité, $f(x)$ varie de $a$.

Sa représentation graphique est une droite.

\section{4. Discret ou continu ?}
Une suite décrit naturellement des étapes entières : mois, années, générations, rangs. Une fonction peut décrire une grandeur sur un intervalle réel. Le choix du modèle dépend donc du phénomène étudié.

\section{5. Problème de seuil}
Pour $u(n)=1400+120n$, cherchons le premier rang tel que $u(n)\ge2000$ :
\[
1400+120n\ge2000
\]
donc
\[
120n\ge600,
\]
d'où $n\ge5$.

Le seuil est atteint au rang 5.

\section{Automatismes à entretenir}
\begin{itemize}
\item Calculer un taux d'évolution.
\item Résoudre une équation ou une inéquation simple du premier degré.
\item Lire le coefficient directeur d'une droite.
\item Tracer une droite à partir de son équation.
\end{itemize}

\section{Erreurs fréquentes}
\begin{itemize}
\item Confondre accroissement constant et taux d'évolution constant.
\item Utiliser un modèle discret pour un phénomène continu sans justification.
\item Oublier que le rang $n$ d'une suite est entier.
\end{itemize}

\section{À retenir}
Croissance linéaire = ajout constant. En discret, le modèle naturel est la suite arithmétique ; en continu, la fonction affine.
